\section{Bezpieczeństwo danych w chmurze obliczeniowej i sformułowanie problemu}
\label{sec:problem}

\subsection{Specyfika zagrożeń w środowisku chmurowym}

Chmura obliczeniowa zmieniła sposób, w jaki firmy przechowują i przetwarzają
dane. Zamiast utrzymywać własne serwerownie, organizacje wynajmują zasoby od
dostawców takich jak Amazon Web Services, Microsoft Azure czy Google Cloud.
Wygoda i oszczędność mają jednak swoją cenę --- dane wielu klientów żyją na tej
samej, współdzielonej infrastrukturze, dostępnej przez sieć z dowolnego miejsca
na świecie. To sprawia, że chmura jest atrakcyjnym celem dla atakujących.

Warto pamiętać o tzw.\ \textbf{modelu współdzielonej odpowiedzialności}: dostawca
chmury dba o bezpieczeństwo samej infrastruktury, ale za poprawną konfigurację,
kontrolę dostępu i ochronę własnych danych odpowiada już klient. W praktyce
ogromna część głośnych wycieków danych z chmury wynikała nie z włamania do
serwerów dostawcy, lecz z błędów po stronie użytkownika --- źle ustawionych
uprawnień, otwartych ,,kubełków'' z danymi czy przejętych kont. Dodatkowo skala
i zmienność ruchu w chmurze są tak duże, że ręczne pilnowanie bezpieczeństwa
przestaje być możliwe.

\subsection{Pojęcie naruszenia danych}

\textbf{Naruszenie danych} (ang.\ \emph{data breach}) to sytuacja, w której
dane trafiają w niepowołane ręce --- zostają wykradzione, ujawnione lub zmienione
bez uprawnienia. Konsekwencje bywają dotkliwe: kary finansowe (np.\ na gruncie
RODO), utrata zaufania klientów i realne straty. Co gorsza, naruszenia często
pozostają niewykryte przez wiele tygodni, co daje atakującym czas na wyrządzenie
dużych szkód.

W środowisku chmurowym najczęstsze drogi prowadzące do naruszenia to m.in.:
\begin{itemize}
  \item \textbf{przejęcie kont} --- np.\ zgadywanie haseł metodą brute force do
        usług logowania;
  \item \textbf{eksfiltracja danych} --- niepostrzeżone wyprowadzanie danych na
        zewnątrz, często ukryte w pozornie zwykłym ruchu;
  \item \textbf{ataki DoS/DDoS} --- zalewanie usług ruchem, czasem jako zasłona
        dymna dla innych działań;
  \item \textbf{infiltracja i ruch boczny} --- po zdobyciu przyczółka atakujący
        przemieszcza się w infrastrukturze ofiary;
  \item \textbf{botnety} oraz \textbf{ataki na aplikacje webowe} (np.\ SQL
        injection), które również prowadzą do wycieku danych.
\end{itemize}

Wszystkie te działania mają jedną wspólną cechę: zostawiają ślad w \emph{ruchu
sieciowym}. Analizując ten ruch, można wychwycić podejrzaną aktywność, zanim
przerodzi się ona w pełne naruszenie.

\subsection{Od wykrywania ataków do prognozowania naruszeń}

Naruszenie danych rzadko jest pojedynczym zdarzeniem --- to zwykle finał całego
łańcucha działań: rozpoznanie, włamanie, eskalacja uprawnień, wreszcie kradzież
danych. Jeśli uda się wykryć wczesne ogniwa tego łańcucha (np.\ nietypowe
połączenie świadczące o skanowaniu lub łamaniu haseł), można zareagować, \emph{zanim}
dojdzie do wycieku. W tym sensie wykrycie ataku pełni rolę \emph{prognostyczną}
--- jest wczesnym sygnałem ostrzegawczym. Dlatego w tej pracy ,,prognozowanie
naruszenia'' utożsamiamy z rozpoznawaniem złośliwych połączeń sieciowych.

\subsection{Sformułowanie zadania klasyfikacji}

Zadanie sprowadza się do prostego pytania zadawanego dla każdego połączenia
sieciowego: \emph{czy to zwykły, normalny ruch, czy raczej atak?} Każde
połączenie opisujemy zestawem liczb --- na przykład jak długo trwało, ile
pakietów i bajtów przesłano, jak często, jakie flagi sieciowe się pojawiły.
Na podstawie tych cech model ma przykleić jedną z dwóch etykiet: \emph{ruch
normalny} albo \emph{atak}. Model uczy się tej umiejętności na historycznych
przykładach, dla których wiadomo, co było czym. Jest to więc klasyczne zadanie
\textbf{klasyfikacji} --- a ponieważ dane zawierają też informację o rodzaju
ataku, możliwa jest dodatkowo analiza poszczególnych typów zagrożeń.

\subsection{Wyzwania badawcze}

Wykrywanie naruszeń nie jest proste z kilku praktycznych powodów:
\begin{itemize}
  \item \textbf{Ataki są rzadkie.} W normalnym ruchu tonie garstka złośliwych
        połączeń. Model łatwo wpada w pułapkę uznawania ,,wszystkiego za
        normalne'' i przeocza zagrożenia.
  \item \textbf{Nie każdy błąd kosztuje tyle samo.} Przeoczenie ataku jest
        zwykle znacznie groźniejsze niż fałszywy alarm --- dlatego zależy nam
        przede wszystkim na wykrywalności (czułości).
  \item \textbf{Złośliwość kryje się w kombinacjach.} Pojedyncza cecha rzadko
        zdradza atak; dopiero ich zestawienie tworzy podejrzany obraz, co
        wyklucza najprostsze, ,,liniowe'' metody.
  \item \textbf{Liczy się szybkość i koszt.} W chmurze model trzeba często
        dotrenowywać i stosować na bieżąco do ogromnych ilości ruchu, więc nie
        może być zbyt kosztowny obliczeniowo.
\end{itemize}

Te właśnie wyzwania kształtują dobór metod i sposób ich oceny, co opisujemy w
kolejnych rozdziałach.
